\input{libraries/report-preamble.tex}
\title{Checked Horner Evaluation on Unsigned Integers}
\begin{document}
\maketitle

\section{Task and interface}
The task is to evaluate $p(x)=\sum_{i=0}^{n-1}c_i x^i$ for nonnegative integer
coefficients and argument.  Each coefficient and $x$ are represented by
\texttt{UInt64}.  The array stores coefficients in ascending degree order.
An empty array represents zero.  Trailing zero coefficients are accepted.

\texttt{LeanExe.Lib.Polynomial.eval} returns \texttt{some value} when every
Horner intermediate fits in a word and \texttt{none} when an intermediate
would exceed $M=2^{64}-1$.  The intended successful result is exact integer
evaluation.  Formal correctness proofs for this component are deferred.

\section{Recurrence and overflow check}
Horner's recurrence processes coefficients from highest degree to lowest,
updating $v\leftarrow vx+c_i$.  The NIST Digital Library of Mathematical
Functions states the recurrence and its evaluation identity~\cite{horner}.
This implementation starts at $v=0$ and processes each stored coefficient.

For $x>0$, the condition $vx+c_i\leq M$ is equivalent to
\[
  v\leq\left\lfloor\frac{M-c_i}{x}\right\rfloor.
\]
The check precedes multiplication and addition.  Subtraction is safe because
$c_i\leq M$.  At $x=0$, each update returns the coefficient, so the check can
be skipped.  The loop reads the borrowed input array and carries an index
and one accumulated value.

\lstinputlisting{LeanExe/Lib/Polynomial/Horner/Basic.lean}

The listing comes from the executable source.  For coefficients $[1,3,2]$
and $x=2$, the accumulated values are $2$, $7$, and $15$.  The implementation
performs $n$ multiply-add updates and at most $n$ divisions for the checks.
This operation count describes the source loop.  Runtime depends on the
compiled program and execution engine.

\section{Clients and commands}
The polynomial client converts the optional result to an explicit status and
value.  The polynomial-ratio client evaluates two arrays, checks the denominator,
and reduces the resulting fraction through the number-theory library.

\begin{lstlisting}
tools/seminum polynomial 2 1,3,2
15
tools/seminum polynomial 99 '[]'
0
tools/seminum ratio 2 1,3,2 2,2
5/2
tools/seminum polynomial 2 18446744073709551615,1
seminum: arithmetic overflow
\end{lstlisting}

The last command writes its diagnostic to standard error and exits with
status 2.  Its exact result would be $M+2$.  Decimal parsing, compilation,
and invocation belong to the command driver.  Polynomial arithmetic executes
in the generated WASM.

\section{Tests and reuse}
The test computes reference values with arbitrary-precision integers using
the sum of powers, compares them with Wasmtime output, and compares the same
cases with native Lean.  It includes empty and constant polynomials, zero
arguments, maximum words, exact boundary results, overflow, and failures
in either polynomial of a ratio.  The example module imports the component
and calls its public declaration.  The compiler's source-name annotations
allow the catalog to identify that declaration in the generated program.

\begin{thebibliography}{9}
\bibitem{horner} NIST Digital Library of Mathematical Functions,
Section 1.11(i), equations 1.11.1--1.11.4.
\url{https://dlmf.nist.gov/1.11.i}.
These equations supply Horner's recurrence and evaluation identity.
The integer overflow check and optional-result API are choices of this
implementation.
\end{thebibliography}
\end{document}
